\documentclass[11pt,a4paper]{article}
\usepackage[utf8]{inputenc}
\usepackage[spanish,es-tabla]{babel}
\usepackage{amsmath,amsfonts,amssymb}
\usepackage{geometry}
\usepackage{graphicx}
\usepackage{xcolor}
\usepackage{enumitem}
\usepackage{tikz}

\geometry{margin=2.5cm}

\title{\textbf{Soluci\'on Examen 2024-2025\\TAF STAR - UEC B\\15/10/2024}}
\author{}
\date{}

\begin{document}

\maketitle

\section*{Preguntas con respuestas cortas (6 puntos)}

\subsection*{Pregunta 1: Principios de dise\~no de constelaci\'on (2 pts)}

\textbf{Enunciado:} Dar los dos principios clave para dise\~nar la constelaci\'on de transmisi\'on de un sistema de comunicaci\'on digital.

\vspace{0.3cm}
\textbf{Soluci\'on:}

Los dos principios fundamentales para el dise\~no \'optimo de una constelaci\'on son:

\textbf{1. Maximizar la distancia euclidiana m\'inima entre s\'imbolos adyacentes:}

\[
\boxed{d_{min} = \min_{i \neq j} |s_i - s_j|}
\]

\textbf{Justificaci\'on:} 

La probabilidad de error en un canal AWGN depende exponencialmente de $d_{min}^2$:
\[
P_e \propto Q\left(\sqrt{\frac{d_{min}^2}{2N_0}}\right)
\]

Maximizar $d_{min}$ minimiza la probabilidad de error.

\textbf{Restricci\'on:} Para una potencia media fija $P_{avg} = \mathbb{E}[|s|^2]$, maximizar $d_{min}$.

\textbf{2. Minimizar el n\'umero de vecinos m\'as cercanos (dada la distancia m\'inima):}

\[
\boxed{N_{min} = \text{n\'umero de s\'imbolos a distancia } d_{min} \text{ de un s\'imbolo dado}}
\]

\textbf{Justificaci\'on:}

La probabilidad de error es aproximadamente:
\[
P_e \approx N_{min} \cdot Q\left(\sqrt{\frac{d_{min}^2}{2N_0}}\right)
\]

Reducir $N_{min}$ reduce proporcionalmente la probabilidad de error.

\textbf{Ejemplos de aplicaci\'on:}

\begin{itemize}
    \item \textbf{BPSK vs On-Off Keying:} BPSK tiene $d_{min}$ el doble que OOK para la misma energ\'ia $\to$ BPSK es \'optima
    
    \item \textbf{16-QAM vs 16-PSK:} Para la misma energ\'ia media:
    \begin{itemize}
        \item 16-QAM: mayor $d_{min}$, s\'imbolos interiores con 4 vecinos
        \item 16-PSK: menor $d_{min}$, todos los s\'imbolos con 2 vecinos
        \item Resultado: 16-QAM es superior (el efecto de $d_{min}$ domina)
    \end{itemize}
    
    \item \textbf{Constelaciones hexagonales:} Reducen $N_{min}$ comparado con ret\'iculas cuadradas
\end{itemize}

\textbf{Orden de prioridad:} El primer principio (maximizar $d_{min}$) es m\'as importante que el segundo, debido al comportamiento exponencial de la funci\'on $Q$.



\subsection*{Pregunta 2: Ancho de banda m\'inimo para transmisi\'on sin ISI (2 pts)}

\textbf{Enunciado:} Consideremos una modulaci\'on 16-QAM. Asumimos que el periodo s\'imbolo es igual a $5 \times 10^{-3}$ segundos y que la frecuencia portadora es igual a 10000 Hertz. Calcular $W_m$ definido como el ancho de banda m\'inimo requerido que asegura transmisi\'on sin ISI.

\vspace{0.3cm}
\textbf{Soluci\'on:}

\textbf{Datos:}
\begin{itemize}
    \item Modulaci\'on: 16-QAM
    \item Periodo s\'imbolo: $T = 5 \times 10^{-3}$ s = 5 ms
    \item Frecuencia portadora: $\nu_0 = 10000$ Hz = 10 kHz
    \item Orden de modulaci\'on: $M = 16$
\end{itemize}

\textbf{Desarrollo:}

Para una modulaci\'on sobre frecuencia portadora (single-carrier modulation), el ancho de banda m\'inimo requerido para transmisi\'on sin ISI est\'a determinado por el \textbf{criterio de Nyquist}.

\textbf{Banda base:}

El ancho de banda m\'inimo en banda base (unilateral) es:
\[
W_{m,baseband} = \frac{1}{2T}
\]

Esto corresponde a un filtro rectangular en frecuencia con soporte $\left[-\frac{1}{2T}, \frac{1}{2T}\right]$.

\textbf{Modulaci\'on sobre portadora:}

Para una modulaci\'on sobre portadora, el espectro se traslada a $\pm\nu_0$, pero el ancho de banda ocupado total permanece igual al de banda base:

\[
\boxed{W_m = \frac{1}{T}}
\]

\textbf{C\'alculo num\'erico:}

\[
W_m = \frac{1}{5 \times 10^{-3}} = 0.2 \times 10^3 = \boxed{200 \text{ Hz}}
\]

\textbf{Explicaci\'on detallada:}

\begin{enumerate}
    \item El filtro de Nyquist ideal (rectangular) tiene banda $\left[-\frac{1}{2T}, \frac{1}{2T}\right]$ en banda base
    
    \item Su respuesta impulsional es $h(t) = \frac{\sin(\pi t/T)}{\pi t/T} = \text{sinc}(t/T)$
    
    \item Este pulso satisface:
    \[
    h(kT) = \begin{cases}
    1 & k = 0 \\
    0 & k \neq 0
    \end{cases}
    \]
    
    \item Garantiza ausencia de ISI con el m\'inimo ancho de banda posible
    
    \item Para la se\~nal modulada:
    \[
    s(t) = \sum_k a_k h(t-kT)\cos(2\pi\nu_0 t) - b_k h(t-kT)\sin(2\pi\nu_0 t)
    \]
    
    El espectro se extiende de $\nu_0 - \frac{1}{2T}$ a $\nu_0 + \frac{1}{2T}$, con ancho total $W_m = \frac{1}{T}$
\end{enumerate}

\textbf{Observaciones importantes:}

\begin{itemize}
    \item El resultado $W_m = 200$ Hz es \textbf{independiente del orden de modulaci\'on} (16-QAM, QPSK, etc.)
    
    \item Solo depende del periodo s\'imbolo $T$, no del d\'ebito binario
    
    \item La frecuencia portadora $\nu_0 = 10$ kHz no afecta el ancho de banda ocupado, solo la posici\'on del espectro
    
    \item En la pr\'actica, se usan filtros con roll-off $\beta > 0$ (ej. raised cosine), que ocupan $W = (1+\beta)/T > W_m$, pero garantizan mejor comportamiento
    
    \item Relaci\'on con la rapidez de s\'imbolo:
    \[
    D_s = \frac{1}{T} = 200 \text{ sym/s} = W_m
    \]
    
    Por tanto: $W_m = D_s$ (en Hz y bauds, respectivamente)
\end{itemize}



\subsection*{Pregunta 3: Elecci\'on de formato de modulaci\'on (2 pts)}

\textbf{Enunciado:} Sea $W = 25000$ Hz la banda ocupada por un sistema de comunicaci\'on. La rapidez de s\'imbolo $D_s$ se impone y $D_s = 12500$ sym/s. El ratio $\frac{E_b}{N_0}$ tambi\'en se impone. Las restricciones de dise\~no del sistema son una probabilidad de error binaria m\'inima a la salida del receptor y una eficiencia espectral m\'inima de 2 bits/s/Hz.

Los formatos de modulaci\'on posibles son 8-PSK, 16-PSK, 16-QAM, 64-QAM y 256-QAM. ¿Cu\'al elegir\'ia para satisfacer las restricciones y por qu\'e?

\vspace{0.3cm}
\textbf{Soluci\'on:}

\textbf{Datos:}
\begin{itemize}
    \item Banda ocupada: $W = 25000$ Hz
    \item Rapidez de s\'imbolo: $D_s = 12500$ sym/s
    \item Ratio $\frac{E_b}{N_0}$ fijo
    \item Restricci\'on: $\eta \geq 2$ bits/s/Hz
    \item Objetivo: Minimizar $P_b$ (probabilidad de error binaria)
\end{itemize}

\textbf{An\'alisis:}

\textbf{Paso 1: Restricci\'on de eficiencia espectral}

La eficiencia espectral se define como:
\[
\eta = \frac{D_b}{W} = \frac{\log_2(M) \cdot D_s}{W}
\]

La restricci\'on $\eta \geq 2$ implica:
\[
\log_2(M) \geq \frac{2W}{D_s} = \frac{2 \times 25000}{12500} = 4
\]

Por lo tanto:
\[
M \geq 2^4 = 16
\]

\textbf{Candidatos que cumplen la restricci\'on:}

\begin{center}
\begin{tabular}{|c|c|c|}
\hline
\textbf{Modulaci\'on} & \textbf{$\log_2(M)$} & \textbf{¿Cumple $\eta \geq 2$?} \\
\hline
8-PSK & 3 & ✗ No ($\eta = 1.5$ bits/s/Hz) \\
16-PSK & 4 & ✓ S\'i ($\eta = 2$ bits/s/Hz) \\
16-QAM & 4 & ✓ S\'i ($\eta = 2$ bits/s/Hz) \\
64-QAM & 6 & ✓ S\'i ($\eta = 3$ bits/s/Hz) \\
256-QAM & 8 & ✓ S\'i ($\eta = 4$ bits/s/Hz) \\
\hline
\end{tabular}
\end{center}

\textbf{Paso 2: An\'alisis de probabilidad de error}

Para un ratio $\frac{E_b}{N_0}$ fijo y diferentes valores de $M$:

\textbf{Comportamiento general:}
\begin{itemize}
    \item Para $M$ fijo: QAM es m\'as eficiente en potencia que PSK (mayor $d_{min}$ para la misma energ\'ia media)
    
    \item Cuando $M$ aumenta: $P_b$ aumenta (m\'as s\'imbolos $\to$ menor distancia entre ellos)
\end{itemize}

\textbf{Comparaci\'on cualitativa (de mejor a peor rendimiento):}

\begin{enumerate}
    \item \textbf{16-QAM:} Mejor compromiso
    \begin{itemize}
        \item $M$ m\'inimo que satisface la restricci\'on
        \item Geometr\'ia QAM (superior a PSK)
        \item Menor $P_b$ para $\frac{E_b}{N_0}$ fijo
    \end{itemize}
    
    \item \textbf{16-PSK:} Rendimiento inferior a 16-QAM
    \begin{itemize}
        \item Mismo $M$ que 16-QAM
        \item Geometr\'ia PSK (inferior a QAM)
        \item Mayor $P_b$ que 16-QAM
    \end{itemize}
    
    \item \textbf{64-QAM:} Peor rendimiento
    \begin{itemize}
        \item $M$ mayor $\to$ menor $d_{min}$
        \item $P_b$ significativamente mayor
    \end{itemize}
    
    \item \textbf{256-QAM:} Mucho peor rendimiento
    \begin{itemize}
        \item $M$ muy grande $\to$ $d_{min}$ muy peque\~no
        \item $P_b$ muy alta
    \end{itemize}
\end{enumerate}

\textbf{Paso 3: Decisi\'on}

\[
\boxed{\text{Modulaci\'on elegida: 16-QAM}}
\]

\textbf{Justificaci\'on:}

\begin{enumerate}
    \item \textbf{Cumple la restricci\'on de eficiencia:}
    \[
    \eta = \frac{4 \times 12500}{25000} = 2 \text{ bits/s/Hz} \quad ✓
    \]
    
    \item \textbf{Minimiza el orden de modulaci\'on:}
    
    Es el $M$ m\'as peque\~no que satisface $\eta \geq 2$, por tanto tiene la mayor distancia m\'inima entre s\'imbolos para la energ\'ia disponible.
    
    \item \textbf{QAM es superior a PSK:}
    
    Para el mismo orden $M = 16$, la constelaci\'on 16-QAM tiene mejor geometr\'ia que 16-PSK:
    \begin{itemize}
        \item Mayor $d_{min}$ para la misma energ\'ia media
        \item Menor probabilidad de error
    \end{itemize}
    
    \item \textbf{Compromiso \'optimo:}
    
    No usar m\'as bits por s\'imbolo de los necesarios, ya que aumentar $M$ degrada el rendimiento.
\end{enumerate}

\textbf{Comparaci\'on cuantitativa aproximada:}

Para SNR moderado, las probabilidades de error relativas son aproximadamente:
\begin{itemize}
    \item 16-QAM: $P_b \approx P_0$ (referencia)
    \item 16-PSK: $P_b \approx 2 P_0$ (el doble)
    \item 64-QAM: $P_b \approx 10 P_0$ (10 veces mayor)
    \item 256-QAM: $P_b \approx 100 P_0$ (100 veces mayor)
\end{itemize}

\textbf{Conclusi\'on:} \boxed{\text{16-QAM es la elecci\'on \'optima}}



\newpage
\section*{Transmisi\'on 8-PAM sobre canal AWGN (14 puntos)}

\textbf{Datos del problema:}
\begin{itemize}
    \item Elementos binarios independientes e id\'enticamente distribuidos: $\Pr(m_i = 0) = \Pr(m_i = 1)$
    \item Modulaci\'on: 8-PAM (Pulse Amplitude Modulation)
\end{itemize}

\textbf{Se\~nal modulada:}
\[
x(t) = \sum_{k=0}^{L-1} a_k h(t-kT)
\]

con
\[
h(t) = \frac{1}{\sqrt{T}} \chi_{[0;T]}(t) = \begin{cases}
\frac{1}{\sqrt{T}} & 0 \leq t \leq T \\
0 & \text{en otro caso}
\end{cases}
\]

donde $\chi_E$ es la funci\'on indicadora del conjunto $E$.

\textbf{Canal:} Gaussiano. Se\~nal recibida: $r(t) = x(t) + w(t)$ con $w(t)$ ruido AWGN de media cero y densidad espectral de potencia $\frac{N_0}{2}$.

\textbf{Receptor:} Filtro con respuesta impulsional $g(t)$ seguido de un muestreador temporal. Sea $z(t)$ la salida del filtro de recepci\'on y $z_\ell$ la $\ell$-\'esima salida del muestreador.

\subsection*{a) Definir el alfabeto elemental 8-PAM (2 pts)}

\textbf{Enunciado:} Definir el alfabeto elemental 8-PAM.

\vspace{0.3cm}
\textbf{Soluci\'on:}

El alfabeto 8-PAM (Pulse Amplitude Modulation) consiste en 8 niveles de amplitud equiespaciados y sim\'etricos:

\[
\boxed{\mathcal{A} = \{-7, -5, -3, -1, +1, +3, +5, +7\}}
\]

\textbf{Propiedades:}
\begin{itemize}
    \item \textbf{N\'umero de s\'imbolos:} $M = 8$
    \item \textbf{Bits por s\'imbolo:} $\log_2(8) = 3$ bits
    \item \textbf{Separaci\'on entre s\'imbolos:} $\Delta = 2$
    \item \textbf{Media:} $\mathbb{E}[a_k] = 0$ (sim\'etrico)
    \item \textbf{Energ\'ia media:}
    \[
    \mathbb{E}[a_k^2] = \frac{1}{8}(49 + 25 + 9 + 1 + 1 + 9 + 25 + 49) = \frac{168}{8} = 21
    \]
    \item \textbf{Distancia m\'inima:} $d_{min} = 2$
\end{itemize}

\textbf{Nota:} PAM es una modulaci\'on en banda base (sin portadora), a diferencia de ASK que modula una portadora.



\subsection*{b) Proponer un Gray mapping (2 pts)}

\textbf{Enunciado:} Proponer un Gray mapping para la conversi\'on binario a s\'imbolo.

\vspace{0.3cm}
\textbf{Soluci\'on:}

\textbf{Gray mapping propuesto:}

\begin{center}
\begin{tabular}{|c|c|c|c|c|}
\hline
\textbf{S\'imbolo $a_k$} & \textbf{$m_{3k}$} & \textbf{$m_{3k+1}$} & \textbf{$m_{3k+2}$} & \textbf{C\'odigo completo} \\
\hline
$-7$ & 0 & 0 & 0 & 000 \\
$-5$ & 0 & 0 & 1 & 001 \\
$-3$ & 0 & 1 & 1 & 011 \\
$-1$ & 0 & 1 & 0 & 010 \\
$+1$ & 1 & 1 & 0 & 110 \\
$+3$ & 1 & 1 & 1 & 111 \\
$+5$ & 1 & 0 & 1 & 101 \\
$+7$ & 1 & 0 & 0 & 100 \\
\hline
\end{tabular}
\end{center}

\textbf{Verificaci\'on del criterio Gray:}

S\'imbolos adyacentes deben diferir en exactamente 1 bit:
\begin{itemize}
    \item $-7 \rightarrow -5$: $000 \rightarrow 001$ (difieren en bit 2) ✓
    \item $-5 \rightarrow -3$: $001 \rightarrow 011$ (difieren en bit 1) ✓
    \item $-3 \rightarrow -1$: $011 \rightarrow 010$ (difieren en bit 2) ✓
    \item $-1 \rightarrow +1$: $010 \rightarrow 110$ (difieren en bit 0) ✓
    \item $+1 \rightarrow +3$: $110 \rightarrow 111$ (difieren en bit 2) ✓
    \item $+3 \rightarrow +5$: $111 \rightarrow 101$ (difieren en bit 1) ✓
    \item $+5 \rightarrow +7$: $101 \rightarrow 100$ (difieren en bit 2) ✓
\end{itemize}

\textbf{Interpretaci\'on de cada bit:}

\begin{itemize}
    \item \textbf{$m_{3k}$ (MSB):} Indica el signo del s\'imbolo
    \begin{itemize}
        \item $m_{3k} = 0 \Rightarrow a_k < 0$ (s\'imbolos negativos)
        \item $m_{3k} = 1 \Rightarrow a_k > 0$ (s\'imbolos positivos)
    \end{itemize}
    
    \item \textbf{$m_{3k+1}$:} Indica la magnitud (peque\~na/grande)
    \begin{itemize}
        \item $m_{3k+1} = 1 \Rightarrow |a_k| \in \{1, 3\}$ (magnitud peque\~na)
        \item $m_{3k+1} = 0 \Rightarrow |a_k| \in \{5, 7\}$ (magnitud grande)
    \end{itemize}
    
    \item \textbf{$m_{3k+2}$ (LSB):} Refina dentro de cada par
\end{itemize}

\textbf{Ventaja:} Este mapping minimiza la tasa de error binaria cuando los errores ocurren entre s\'imbolos adyacentes (los m\'as probables en AWGN).



\subsection*{d) Demostrar que $h(t)$ satisface TSWOC (3 pts)}

\textbf{Enunciado:} Demostrar que el filtro de transmisi\'on (definido por $h(t) = \frac{1}{\sqrt{T}}\chi_{[0;T]}(t)$) satisface el criterio TSWOC. Se espera una expresi\'on cuidadosa del criterio y la demostraci\'on puede basarse en una representaci\'on gr\'afica relevante.

\vspace{0.3cm}
\textbf{Soluci\'on:}

\textbf{Criterio TSWOC (Time-Shifted Waveform Orthogonality Criterion):}

\[
\boxed{\int_{\mathbb{R}} h(t) h(t - kT)\, dt = \delta_{k,0}}
\]

donde $\delta_{k,0}$ es la delta de Kronecker:
\[
\delta_{k,0} = \begin{cases}
1 & \text{si } k = 0 \\
0 & \text{si } k \neq 0
\end{cases}
\]

\textbf{Demostraci\'on:}

Calculamos la integral para los diferentes valores de $k$:

\textbf{Caso 1: $k = 0$}

\begin{align*}
\int_{\mathbb{R}} h(t) h(t)\, dt &= \int_{\mathbb{R}} |h(t)|^2\, dt \\
&= \int_0^T \left(\frac{1}{\sqrt{T}}\right)^2\, dt \\
&= \frac{1}{T} \int_0^T dt \\
&= \frac{1}{T} \cdot T \\
&= 1 \quad ✓
\end{align*}

\textbf{Caso 2: $k \neq 0$}

Necesitamos calcular:
\[
\int_{\mathbb{R}} h(t) h(t - kT)\, dt = \int_{\mathbb{R}} \frac{1}{\sqrt{T}}\chi_{[0;T]}(t) \cdot \frac{1}{\sqrt{T}}\chi_{[kT;(k+1)T]}(t)\, dt
\]

El producto $\chi_{[0;T]}(t) \cdot \chi_{[kT;(k+1)T]}(t)$ solo es no nulo donde ambos intervalos se superponen.

\textbf{An\'alisis de intersecci\'on:}
\begin{itemize}
    \item $h(t)$ tiene soporte en $[0, T]$
    \item $h(t - kT)$ tiene soporte en $[kT, (k+1)T]$
\end{itemize}

Para $k \geq 1$: $[0, T] \cap [kT, (k+1)T] = \emptyset$ (intervalos disjuntos)

Para $k \leq -1$: $[0, T] \cap [kT, (k+1)T] = \emptyset$ (intervalos disjuntos)

Por lo tanto:
\[
\int_{\mathbb{R}} h(t) h(t - kT)\, dt = \frac{1}{T} \int_{\emptyset} dt = 0 \quad ✓
\]

\textbf{Representaci\'on gr\'afica:}

\begin{center}
\begin{tikzpicture}[scale=1.2]
% Eje temporal
\draw[->] (-1,0) -- (5,0) node[right] {$t$};
\draw[->] (0,-0.5) -- (0,2.5) node[above] {$h(t)$};

% h(t)
\draw[blue, very thick] (0,0) -- (0,2) -- (2,2) -- (2,0);
\node[blue, above] at (1,2) {$h(t)$};
\node[below] at (0,-0.1) {$0$};
\node[below] at (2,-0.1) {$T$};

% h(t - T) (k=1)
\draw[red, dashed, thick] (2,0) -- (2,1.5) -- (4,1.5) -- (4,0);
\node[red, above] at (3,1.5) {$h(t-T)$};
\node[below] at (4,-0.1) {$2T$};

% Indicaci\'on de no solapamiento
\draw[<->, thick, green!70!black] (0,-0.3) -- (2,-0.3);
\node[green!70!black, below] at (1,-0.3) {Soporte de $h(t)$};

\draw[<->, thick, orange!70!black] (2,-0.6) -- (4,-0.6);
\node[orange!70!black, below] at (3,-0.6) {Soporte de $h(t-T)$};

\node at (2.5, 2.3) {\textbf{No hay solapamiento} $\to$ Integral = 0};

\end{tikzpicture}
\end{center}

\textbf{Conclusi\'on:}

\[
\boxed{\int_{\mathbb{R}} h(t) h(t - kT)\, dt = \begin{cases}
1 & k = 0 \\
0 & k \neq 0
\end{cases} = \delta_{k,0}}
\]

Por lo tanto, \textbf{el filtro rectangular $h(t) = \frac{1}{\sqrt{T}}\chi_{[0;T]}(t)$ satisface el criterio TSWOC}.

\textbf{Interpretaci\'on f\'isica:}

Este resultado significa que pulsos rectangulares sucesivos son ortogonales entre s\'i cuando est\'an desplazados por m\'ultiplos enteros del periodo $T$. Esta propiedad garantiza:
\begin{itemize}
    \item Ausencia de interferencia entre s\'imbolos (ISI)
    \item Detecci\'on \'optima s\'imbolo por s\'imbolo
    \item M\'aximo SNR en el receptor con filtro adaptado
\end{itemize}



\subsection*{e) Expresi\'on de $g(t)$ para detecci\'on ML \'optima (1 pt)}

\textbf{Enunciado:} Consideramos primero que $t_0$ est\'a fijado. Dados los par\'ametros de modulaci\'on, las hip\'otesis y la estructura del receptor, escribir la expresi\'on de $g(t)$ que asegura detecci\'on ML \'optima s\'imbolo por s\'imbolo.

\vspace{0.3cm}
\textbf{Soluci\'on:}

Para un receptor \'optimo de m\'axima verosimilitud (ML) que garantice detecci\'on s\'imbolo por s\'imbolo, el filtro de recepci\'on debe ser el \textbf{filtro adaptado} al filtro de transmisi\'on:

\[
\boxed{g(t) = h(t_0 - t)}
\]

Sustituyendo la expresi\'on de $h(t)$:

\[
h(t_0 - t) = \frac{1}{\sqrt{T}} \chi_{[0;T]}(t_0 - t)
\]

Para que $(t_0 - t) \in [0, T]$, necesitamos $t \in [t_0 - T, t_0]$.

Por lo tanto:

\[
\boxed{g(t) = \frac{1}{\sqrt{T}} \chi_{[t_0-T; t_0]}(t) = \begin{cases}
\frac{1}{\sqrt{T}} & t_0 - T \leq t \leq t_0 \\
0 & \text{en otro caso}
\end{cases}}
\]

\textbf{Interpretaci\'on:}

\begin{itemize}
    \item $g(t)$ es una versi\'on ''volteada'' y desplazada de $h(t)$
    \item Si $h(t)$ es un pulso rectangular en $[0, T]$, entonces $g(t)$ es un pulso rectangular en $[t_0-T, t_0]$
    \item $t_0$ debe ser al menos $T$ para garantizar causalidad: $t_0 \geq T$
\end{itemize}

\textbf{Elecci\'on t\'ipica:} $t_0 = T$, lo que da:
\[
g(t) = \frac{1}{\sqrt{T}} \chi_{[0;T]}(t) = h(t)
\]

En este caso especial, el filtro de recepci\'on es id\'entico al filtro de transmisi\'on (el pulso rectangular es sim\'etrico).

\textbf{Propiedades del filtro adaptado:}

\begin{enumerate}
    \item \textbf{Maximiza el SNR:} En el instante de muestreo $t_0$, el SNR es m\'aximo
    
    \item \textbf{Combinado con TSWOC:} Garantiza ausencia de ISI
    
    \item \textbf{Implementaci\'on:} $(h * g)(t) = (h * h)(t_0 - t)$ es la autocorrelaci\'on de $h$
\end{enumerate}



\subsection*{f) Expresi\'on de $z_\ell$ (2 pts)}

\textbf{Enunciado:} Asumiendo condiciones satisfechas para detecci\'on ML \'optima s\'imbolo por s\'imbolo, deducir la expresi\'on de $z_\ell$.

\vspace{0.3cm}
\textbf{Soluci\'on:}

\textbf{Desarrollo:}

La salida del filtro de recepci\'on es:
\[
z(t) = (r * g)(t) = \int_{\mathbb{R}} r(\tau) g(t - \tau)\, d\tau
\]

Sustituyendo $r(t) = x(t) + w(t)$ y $x(t) = \sum_{k=0}^{L-1} a_k h(t-kT)$:

\begin{align*}
z(t) &= \int_{\mathbb{R}} \left[\sum_{k=0}^{L-1} a_k h(\tau-kT) + w(\tau)\right] g(t-\tau)\, d\tau \\
&= \sum_{k=0}^{L-1} a_k \int_{\mathbb{R}} h(\tau-kT) g(t-\tau)\, d\tau + \int_{\mathbb{R}} w(\tau) g(t-\tau)\, d\tau \\
&= \sum_{k=0}^{L-1} a_k (h * g)(t-kT) + (w * g)(t)
\end{align*}

Con el cambio de variable $\tau' = \tau - kT$:
\[
(h * g)(t-kT) = \int_{\mathbb{R}} h(\tau') g(t - kT - \tau')\, d\tau' = \int_{\mathbb{R}} h(\tau') g((t-kT) - \tau')\, d\tau'
\]

Como $g(t) = h(t_0 - t)$:
\[
(h * g)(t-kT) = \int_{\mathbb{R}} h(\tau') h(t_0 - (t - kT) + \tau')\, d\tau' = \int_{\mathbb{R}} h(\tau') h(\tau' + t_0 - t + kT)\, d\tau'
\]

\textbf{Muestreo en $t = t_0 + \ell T$:}

\begin{align*}
z_\ell &= z(t_0 + \ell T) \\
&= \sum_{k=0}^{L-1} a_k \int_{\mathbb{R}} h(\tau) h(\tau + t_0 - (t_0 + \ell T) + kT)\, d\tau + w_\ell \\
&= \sum_{k=0}^{L-1} a_k \int_{\mathbb{R}} h(\tau) h(\tau + (k - \ell)T)\, d\tau + w_\ell
\end{align*}

donde $w_\ell = (w * g)(t_0 + \ell T)$ es el ruido filtrado y muestreado.

\textbf{Aplicaci\'on del TSWOC:}

Como demostramos que $\int_{\mathbb{R}} h(\tau) h(\tau + mT)\, d\tau = \delta_{m,0}$, solo el t\'ermino con $k = \ell$ sobrevive:

\[
z_\ell = a_\ell \int_{\mathbb{R}} |h(\tau)|^2\, d\tau + w_\ell = a_\ell \cdot 1 + w_\ell
\]

Por lo tanto:

\[
\boxed{z_\ell = a_\ell + w_\ell}
\]

Donde:
\begin{itemize}
    \item $a_\ell \in \{-7, -5, -3, -1, +1, +3, +5, +7\}$ es el s\'imbolo transmitido
    \item $w_\ell$ es una variable aleatoria gaussiana con:
    \begin{itemize}
        \item Media: $\mathbb{E}[w_\ell] = 0$
        \item Varianza: $\sigma_w^2 = \frac{N_0}{2}$
    \end{itemize}
\end{itemize}

\textbf{Interpretaci\'on:}

La muestra $z_\ell$ es simplemente el s\'imbolo transmitido $a_\ell$ m\'as ruido aditivo gaussiano. No hay interferencia de otros s\'imbolos (ISI = 0) gracias al TSWOC.

Esta es la situaci\'on ideal para detecci\'on, permitiendo decisi\'on s\'imbolo por s\'imbolo con probabilidad de error m\'inima.



\subsection*{g) Constelaci\'on recibida y regiones de decisi\'on (2 pts)}

\textbf{Enunciado:} Dibujar la constelaci\'on recibida. Definir y dibujar las fronteras de decisi\'on sobre los s\'imbolos (con una frase de justificaci\'on). A\~nadir el mapping binario a s\'imbolo propuesto previamente.

\vspace{0.3cm}
\textbf{Soluci\'on:}

\textbf{Constelaci\'on recibida (sin ruido):}

Con ruido cero ($w_\ell = 0$), la constelaci\'on recibida es simplemente:
\[
\{-7, -5, -3, -1, +1, +3, +5, +7\}
\]

\textbf{Definici\'on de fronteras de decisi\'on:}

Las fronteras de decisi\'on se ubican \textbf{a mitad de camino entre s\'imbolos adyacentes} para minimizar la probabilidad de error (criterio ML).

Fronteras: $-6, -4, -2, 0, +2, +4, +6$

\textbf{Representaci\'on gr\'afica:}

\begin{center}
\begin{tikzpicture}[scale=0.9]
% Eje
\draw[<->] (-8.5,0) -- (8.5,0) node[right] {$z_\ell$};

% S\'imbolos
\foreach \x/\label in {-7/000, -5/001, -3/011, -1/010, 1/110, 3/111, 5/101, 7/100}
{
    \fill[blue] (\x,0) circle (4pt);
    \node[blue, above] at (\x,0.5) {\textbf{\label}};
    \node[below] at (\x,-0.3) {$\x$};
}

% Fronteras de decisi\'on
\foreach \x in {-6, -4, -2, 0, 2, 4, 6}
{
    \draw[red, thick, dashed] (\x,-1) -- (\x,1);
    \node[red, below] at (\x,-1.2) {$\x$};
}

% Regiones de decisi\'on
\node[green!70!black] at (-7,-1.8) {Regi\'on 000};
\node[green!70!black] at (-5,-1.8) {001};
\node[green!70!black] at (-3,-1.8) {011};
\node[green!70!black] at (-1,-1.8) {010};
\node[green!70!black] at (1,-1.8) {110};
\node[green!70!black] at (3,-1.8) {111};
\node[green!70!black] at (5,-1.8) {101};
\node[green!70!black] at (7,-1.8) {100};

\end{tikzpicture}
\end{center}

\textbf{Justificaci\'on de las fronteras:}

\textit{Las fronteras de decisi\'on se colocan en los puntos equidistantes entre s\'imbolos adyacentes para minimizar la probabilidad de error en un canal AWGN. Cualquier valor recibido $z_\ell$ se asigna al s\'imbolo m\'as cercano.}

\textbf{Reglas de decisi\'on:}

\[
\hat{a}_\ell = \begin{cases}
-7 & \text{si } z_\ell < -6 \\
-5 & \text{si } -6 \leq z_\ell < -4 \\
-3 & \text{si } -4 \leq z_\ell < -2 \\
-1 & \text{si } -2 \leq z_\ell < 0 \\
+1 & \text{si } 0 \leq z_\ell < 2 \\
+3 & \text{si } 2 \leq z_\ell < 4 \\
+5 & \text{si } 4 \leq z_\ell < 6 \\
+7 & \text{si } z_\ell \geq 6
\end{cases}
\]

\textbf{Reglas de decisi\'on binarias:}

Basadas en el Gray mapping propuesto:

\textbf{Bit $\hat{m}_{3\ell}$ (MSB - Signo):}
\[
\boxed{\hat{m}_{3\ell} = \begin{cases}
0 & \text{si } z_\ell < 0 \\
1 & \text{si } z_\ell \geq 0
\end{cases}}
\]

\textbf{Bit $\hat{m}_{3\ell+1}$ (Magnitud):}
\[
\boxed{\hat{m}_{3\ell+1} = \begin{cases}
1 & \text{si } |z_\ell| < 4 \\
0 & \text{si } |z_\ell| \geq 4
\end{cases}}
\]

\textbf{Bit $\hat{m}_{3\ell+2}$ (LSB - Refinamiento):}
\[
\boxed{\hat{m}_{3\ell+2} = \begin{cases}
1 & \text{si } 2 < |z_\ell| < 6 \\
0 & \text{si } |z_\ell| < 2 \text{ o } |z_\ell| > 6
\end{cases}}
\]



\subsection*{h) Bonus: Reglas de decisi\'on binarias (1 pt)}

\textbf{Enunciado:} Deducir del esquema las reglas de decisi\'on binarias.

\vspace{0.3cm}
\textbf{Soluci\'on:}

Ya desarrolladas en el apartado anterior. Resumen:

\textbf{Implementaci\'on pr\'actica:}

\begin{enumerate}
    \item \textbf{Para $\hat{m}_{3\ell}$:} Comparador simple
    \[
    \hat{m}_{3\ell} = \text{signo}(z_\ell) = \begin{cases}
    0 & z_\ell < 0 \\
    1 & z_\ell \geq 0
    \end{cases}
    \]
    
    \item \textbf{Para $\hat{m}_{3\ell+1}$:} Comparador de magnitud
    \[
    \hat{m}_{3\ell+1} = \begin{cases}
    1 & |z_\ell| < 4 \\
    0 & |z_\ell| \geq 4
    \end{cases}
    \]
    
    \item \textbf{Para $\hat{m}_{3\ell+2}$:} Dos comparadores
    \[
    \hat{m}_{3\ell+2} = \mathbb{1}_{2 < |z_\ell| < 6} = \begin{cases}
    1 & 2 < |z_\ell| < 6 \\
    0 & \text{en otro caso}
    \end{cases}
    \]
\end{enumerate}

\textbf{Circuito decisor:}

\begin{center}
\begin{tikzpicture}[node distance=2cm]
\node (input) {$z_\ell$};
\node[draw, right of=input] (abs) {$|z_\ell|$};
\node[draw, above right of=abs] (comp1) {$< 0?$};
\node[draw, right of=abs] (comp2) {$< 4?$};
\node[draw, below right of=abs] (comp3) {$2 < \cdot < 6?$};

\draw[->] (input) -- (abs);
\draw[->] (input) |- (comp1);
\draw[->] (abs) -- (comp2);
\draw[->] (abs) -- (comp3);

\node[right of=comp1] (out1) {$\hat{m}_{3\ell}$};
\node[right of=comp2] (out2) {$\hat{m}_{3\ell+1}$};
\node[right of=comp3] (out3) {$\hat{m}_{3\ell+2}$};

\draw[->] (comp1) -- (out1);
\draw[->] (comp2) -- (out2);
\draw[->] (comp3) -- (out3);
\end{tikzpicture}
\end{center}

\textbf{Ventaja del Gray mapping:}

Los errores m\'as probables (confusi\'on con s\'imbolo adyacente) causan solo 1 bit err\'oneo, minimizando la tasa de error binaria.



\vspace{1cm}
\begin{center}
\rule{0.8\linewidth}{0.4pt}\\
\textit{Fin del examen}
\end{center}

\end{document}
